\section{Exercises}\label{sec:07-groups:07-exercises:1}

\textbf{Key points.} A group is data (\texttt{op}, \texttt{id}, \texttt{inv}) plus proof obligations (\texttt{assoc}, two-sided \texttt{id}/\texttt{inv} laws), bundled in one \texttt{structure}. The left/right split on \texttt{id}/\texttt{inv} is not overcaution: \texttt{perm3Group} needs both directions, genuinely different. Any theorem proved once, generically, against \texttt{Grp : Group G} applies to every concrete group built afterward at no extra cost.

\begin{enumerate}
\def\labelenumi{\arabic{enumi}.}
\item
  \texttt{GroupData} (\texttt{op}, \texttt{id}, \texttt{inv}, no axioms) type-checks for any choice of the three fields, including nonsense ones. State which fields \texttt{Group} adds to cut out the genuine groups. Then exhibit a \texttt{GroupData Bool} whose operation is associative and has a left identity, but no right identity, and prove it, showing that \texttt{id\_\allowbreak{}right} cannot be derived from the other four axioms and must remain a field of its own.
\item
  Build \texttt{boolXorGroup : Group Bool} where \texttt{op} is boolean XOR (\texttt{Bool.xor}), \texttt{id := false}, and \texttt{inv := fun a => a} (every element is its own inverse). \texttt{by intro a; cases a <;> rfl} proves each field in one line; instead use \texttt{cases a with | false => rfl | true => rfl} for the fields that need a case split, to see which case does what.
\item
  Prove: if \texttt{Grp : Group G} and \texttt{Grp.op} is commutative, then \texttt{id\_\allowbreak{}left} and \texttt{id\_\allowbreak{}right} are logically equivalent, and likewise \texttt{inv\_\allowbreak{}left} and \texttt{inv\_\allowbreak{}right}. Then exhibit, using \texttt{perm3Group}, why the two directions cannot be collapsed into one axiom in general.
\item
  State the general principle that lets a single computed inequality \(f(x) \neq g(x)\) at one point \(x\) serve as a complete proof that \(f \neq g\) as functions, and prove it. Apply it to the two \texttt{\#eval}s in Section 4 to conclude \texttt{perm3Group} is non-abelian.
\end{enumerate}

Solutions, \hyperref[sec:15-appendix-solutions:07-chapter-7:1]{Appendix, Chapter 7}.

